%!TEX root = ../template.tex
%%%%%%%%%%%%%%%%%%%%%%%%%%%%%%%%%%%%%%%%%%%%%%%%%%%%%%%%%%%%%%%%%%%
%% chapter1.tex
%% NOVA thesis document file
%%
%% Chapter with introduction
%%%%%%%%%%%%%%%%%%%%%%%%%%%%%%%%%%%%%%%%%%%%%%%%%%%%%%%%%%%%%%%%%%%

\typeout{NT FILE chapter1.tex}%

\chapter{Introduction}
\label{cha:introduction}

\prependtographicspath{{5-Figures/Covers/}}

% epigraph configuration
\epigraphfontsize{\small\itshape}
\setlength\epigraphwidth{12.5cm}
\setlength\epigraphrule{0pt}

% \section{Welcome to the \novathesis\ Template}
% \label{sec:if_you_use_this_template}

\paragraph{}As the field of telecommunications progresses issues such as scalability, routing, quality of
service, resource management and network security become increasingly complex. Recent studies 
turn to Artificial Intelligence (AI) and Machine Learning (ML) for solutions, specifically MADDPG and GNN models involving DRL 
agents to make optimal routing decisions.

The convergence of Deep Reinforcement Learning (DRL), Graph Neural Networks (GNNs), 
and automatization of network management and configuration represents a radical change in how modern networks 
are controlled and optimized. Traditional routing protocols, such as OSPF and BGP, 
employ static rule-based approaches that struggle with dynamic network environments characterized by heterogeneous Quality-of-Service (QoS) 
requirements, variable traffic patterns, and multiple optimization constraints. 
Deep Reinforcement Learning has emerged as a powerful alternative, enabling routing agents 
to learn adaptive policies directly from real-time network observations without explicit 
programming of optimization rules. \cite{Tian2023maddpgslice}

Multi-Agent DRL (MADRL), more specifically utilizing Multi-Agent Deep Deterministic Policy Gradient (MADDPG), systems extend this capability to distributed network environments 
where multiple autonomous agents must coordinate routing decisions across different nodes 
while maintaining local operational autonomy. When enhanced with Graph Neural Networks that 
naturally exploit the graph-structured topology of networks, these systems achieve improved 
scalability, topology-awareness, and the ability to generalize across different network 
configurations. \cite{wang2024multi} \cite{huang2025era} \cite{9974607} \cite{zhang2025dramadynamicpacketrouting} \cite{11134361}.

\paragraph{}However DRL models have many well documented security vulnerabilities, consequences of which
remain underexplored in the context of network routing. Common attacks such as the fast gradient sign method
(FGSM) \cite{goodfellow2015explainingharnessingadversarialexamples}, the projected gradient descent (PGD)
\cite{madry2019deeplearningmodelsresistant} and the Carlini and Wagner (CW) attack 
\cite{carlini2017evaluatingrobustnessneuralnetworks} can manipulate model inputs to degrade performance or even implant backdoors
as a gateway to further exploits \cite{dai2025effectivebackdoorattackgraph}. Some research has been done on the effects of adversarial 
attacks on DRL models in the context of image classification, mostly focusing on 
model robustness versus adversarial examples. \cite{wang2022attackingdefendingdeepreinforcement}
\cite{goodfellow2015explainingharnessingadversarialexamples} \cite{villegas2024evaluating}.

\paragraph{}This thesis aims to explore the effects of adversarial attacks on MADDPG and GNN models in the context of network routing.
Specifically, we will investigate how these attacks impact the performance of these models when used for routing
decisions in a simulated network environment. We will implement and adapt FGSM and PGD algorithms to the network context,
evaluating their effectiveness in degrading network performance metrics accross diverse scenarios.

\paragraph{}The main contributions of this thesis are as follows:
\begin{itemize}
    \item Adapting and validating FGSM and PGD adversarial attack algorithms to the 
    context of network routing.
    \item Evaluating the impact of these attacks on the performance of MADDPG and 
    GNN routing agents across various network topologies and traffic patterns.
\end{itemize}

\paragraph{} The structure of this thesis preparation report is as follows: Chapter 1 provides an introduction
to the research topic, outlining the motivation and objectives of the study. Chapter 2 reviews relevant literature,
summarizing existing research on MADDPG, GNNs and adversarial attacks of respective architectures.
Chapter 3 details the methodology and implementation of the working enviroment, FGSM and PGD adapted to the netwrok enviroment.